%@ pages 395-402
%@ status translated
% Chapter IV, section 4: Duality between locally finite categories and
% pseudocompact modules (pp. 395--402). Transcribed from the page images
% in reference/source.pdf, not from the OCR text layer, which mangles
% subscripts, limits and fraktur letters. The section opens with its
% heading in the middle of p. 395 and runs to the middle of p. 402
% ("Ceci achève la preuve." at the end of the proof of Corollary 6),
% where Chapter V begins.
%
% Numbering is automatic (src/preamble.tex). This section prints
% Propositions 12 and 13 (propositions run continuously across the sections
% of the chapter; section~3 ends with Proposition 11), Theorem 4 (theorems
% restart at each chapter), Lemmas 1 and 2 (the lemma counter restarts at
% every section) and Corollaries 1--6 (the corollary counter restarts at
% every section).
%
% No emendations.

\section{Duality between Locally Finite Categories and Pseudocompact Modules}

--- We propose to show that every locally finite category is equivalent
to the dual category of a category of pseudocompact modules; we begin
with a few definitions:

A full subcategory of a category $\mathbf{A}$ is said to be
\emph{closed} if the following conditions are satisfied: every
subobject and every quotient of an object of $\mathbf{C}$ belongs to
$\mathbf{C}$; the direct sum (in $\mathbf{A}$) of two objects of
$\mathbf{C}$ belongs to $\mathbf{C}$; the inductive limit (in
$\mathbf{A}$) of an inductive system of objects of $\mathbf{C}$ belongs
to $\mathbf{C}$.

If $M$ is an object of $\mathbf{A}$, $M$ contains a subobject
$\mathbf{C}M$ which belongs to $\mathbf{C}$ and which contains all the
other subobjects of $M$ belonging to $\mathbf{C}$. If $f\colon M\to N$
is a morphism of $\mathbf{A}$, $f(\mathbf{C}M)$ is contained in
$\mathbf{C}N$; moreover, the functor $M\mapsto\mathbf{C}M$ is adjoint
to the canonical functor from $\mathbf{C}$ into $\mathbf{A}$; in
particular, the functor $M\mapsto\mathbf{C}M$ is left exact
(Proposition~11, Chap.~I). For example, the semisimple objects of
$\mathbf{A}$ are the objects of a closed subcategory $\mathbf{S}$. If
$M$ is an object of $\mathbf{A}$, the subobject $\mathbf{S}M$ is called
the \emph{socle} of $M$.

If $\mathbf{C}$ and $\mathbf{D}$ are two closed subcategories of
$\mathbf{A}$, we denote by $\mathbf{C}\cdot\mathbf{D}$ the following
closed subcategory: an object $M$ of $\mathbf{A}$ belongs to
$\mathbf{C}\cdot\mathbf{D}$ if $M/\mathbf{D}M$ is an object of
$\mathbf{C}$. The product
$(\mathbf{C},\mathbf{D})\mapsto\mathbf{C}\cdot\mathbf{D}$ is manifestly
associative. We may therefore speak of the powers $\mathbf{C}^n$ of a
closed subcategory $\mathbf{C}$ of $\mathbf{A}$.

Let now $\mathbf{A}$ be a locally finite category and $\mathbf{S}$ the
closed subcategory of $\mathbf{A}$ defined by the semisimple objects of
$\mathbf{A}$. If $M$ is an object of $\mathbf{A}$, we denote by $M_n$
the largest subobject of $M$ belonging to the closed subcategory
$\mathbf{S}^{n+1}$ (for example, $M_0$ will denote the socle of $M$).
Since $\mathbf{A}$ is a locally finite category, it is clear that $M$
is the supremum of the subobjects $M_n$. We shall examine the case
where $M$ is an injective object of $\mathbf{A}$.

\begin{proposition}
Let $I$ be an injective object of a locally finite category
$\mathbf{A}$, $A$ the ring of endomorphisms of $I$ and $\mathfrak{r}$
the set of endomorphisms of $I$ whose kernel contains the socle $I_0$.
Then $\mathfrak{r}$ is the Jacobson radical of $A$; the intersection of
the powers $\mathfrak{r}^n$ is zero; the quotient ring
$A/\mathfrak{r}$ is isomorphic to the product of the rings of
endomorphisms of a family of vector spaces.
\end{proposition}

Let indeed $\mathfrak{r}^{(n)}$ be the two-sided ideal formed of the
endomorphisms of $I$ whose kernel contains $I_{n-1}$. If $f$ is an
element of $\mathfrak{r}^{(n)}$, it is clear that $f(I_{m+n})$ is
contained in $I_m$. It follows that
$\mathfrak{r}^{(m)}\cdot\mathfrak{r}^{(n)}$ is contained in
$\mathfrak{r}^{(m+n)}$. In particular $\mathfrak{r}^n$ is contained in
$\mathfrak{r}^{(n)}$.

On the other hand the canonical morphisms from $I_n$ into $I$ and from
$I$ onto $I/I_n$ give rise to exact sequences
\[
0\to\operatorname{Hom}(I/I_n,I)\xrightarrow{i_n}\operatorname{Hom}(I,I)\xrightarrow{p_n}\operatorname{Hom}(I_n,I)\to0.
\]

The ideal $\mathfrak{r}^{(n+1)}$ is none other than the image of $i_n$.
The map $p_n$ therefore defines an isomorphism from
$A/\mathfrak{r}^{(n+1)}$ onto $\operatorname{Hom}(I_n,I)$. The formulas
\[
A=\operatorname{Hom}(I,I)=\operatorname{Hom}(\varinjlim I_n,I)=\varprojlim\operatorname{Hom}(I_n,I)=\varprojlim A/\mathfrak{r}^{(n+1)}
\]
then show that $A$ is separated and complete for the filtration defined
by the ideals $\mathfrak{r}^{(n)}$. The formula
$(1-x)^{-1}=1+x+x^2+\cdots$ shows moreover that the inverse of $1-x$
exists if $x$ belongs to $\mathfrak{r}$; in other words,
$\mathfrak{r}$ is contained in the Jacobson radical of $A$.

Let us note finally that for every morphism $f$ from $I_0$ into $I$,
$f(I_0)$ is contained in $I_0$. The quotient $A/\mathfrak{r}$ is
therefore equal to $\operatorname{Hom}(I_0,I_0)$: it is the ring of
endomorphisms of a semisimple object. The last assertion of the
proposition follows from this (cf.\ Chap.~I, \S~6). It also follows
that $\mathfrak{r}$ is an intersection of maximal ideals and contains
the Jacobson radical; this completes the proof.

When $M$ runs through the subobjects of finite length of $I$, the
canonical morphism from $I$ onto $I/M$ defines an injection of
$\operatorname{Hom}(I/M,I)$ into $A$. The image of this map is a left
ideal which we shall denote $l(M)$. If $M$ and $N$ are two subobjects
of finite length, the exactness of the functor
$X\mapsto\operatorname{Hom}(X,I)$ implies the equality
$l(M+N)=l(M)\cap l(N)$. This equality shows that the left ideals $l(M)$
form a base of neighbourhoods of $0$ for a topology which makes $A$ a
topological group. We shall always equip $A$ with this topology, which
we shall call \emph{natural}.

\begin{proposition}
a.\ Let $\mathbf{A}$ be a locally finite category, $I$ an injective of
$\mathbf{A}$ and $A$ the ring of endomorphisms of $I$. Equipped with the
natural topology, $A$ is a pseudocompact ring.

b.\ Conversely, every pseudocompact ring $A$ is isomorphic to the ring
of endomorphisms of an injective of a locally finite category. In
particular, the Jacobson radical $\mathfrak{r}$ of $A$ is the
intersection of the maximal open left ideals; the intersection of the
$\mathfrak{r}^n$ is zero and $A/\mathfrak{r}$ is isomorphic to the
product of the rings of endomorphisms of a family of vector spaces.
\end{proposition}

a.\ Let us first show that $A$ is a topological ring. Since there
exists a base of neighbourhoods of $0$ formed of left ideals, it
suffices to prove the following: for every subobject of finite length
$M$ of $I$ and for every element $a$ of $A$, there exists a subobject
of finite length $N$ such that $l(N)\cdot a$ is contained in $l(M)$.
But it suffices to choose $N$ equal to $a(M)$.

We have on the other hand the ``equalities''
\[
A=\operatorname{Hom}(I,I)=\operatorname{Hom}(\varinjlim M,I)=\varprojlim\operatorname{Hom}(M,I)=\varprojlim A/l(M).
\]
These equalities show both that $A$ is separated and that $A$ is
complete. It remains to prove that the $A$-modules $A/l(M)$ are of
finite length: for this, we first study the $A$-module
$\operatorname{Hom}(S,I)$ when $S$ is a simple object of $\mathbf{A}$;
we denote by $I_S$ the isotypic component of type $S$ of $I_0$ and by
$B$ the ring of endomorphisms of $I_S$.

For every endomorphism $f$ of $I$, $f(I_S)$ is contained in $I_S$; it
follows that $f$ induces an endomorphism $g$ of $I_S$ and that the map
$p\colon f\mapsto g$ from $A$ into $B$ is an epimorphism of rings. This
epimorphism is compatible with the canonical bijection of the
$A$-module $\operatorname{Hom}(S,I)$ onto the $B$-module
$\operatorname{Hom}(S,I_S)$. It follows that $\operatorname{Hom}(S,I)$
is zero or is a simple $A$-module according as $I_S$ is zero or not.

Let now $M$ be a subobject of finite length of $I$ and let
\[
0\subset M_1\subset M_2\subset\cdots\subset M_n\subset M
\]
be a Jordan--Hölder series of $M$. The $A$-module $A/l(M)$ is
isomorphic to $\operatorname{Hom}(M,I)$ and the modules
$\operatorname{Hom}(M/M_i,I)$ define a composition series of
$\operatorname{Hom}(M,I)$. The quotient of $\operatorname{Hom}(M/M_i,I)$
by $\operatorname{Hom}(M/M_{i+1},I)$ is isomorphic to
$\operatorname{Hom}(M_{i+1}/M_i,I)$. It follows that the length of
$A/l(M)$ is less than or equal to the length of $M$. This proves (a).

b.\ Let $\mathbf{A}$ be the dual category of $\operatorname{PC}(A)$.
The left $A$-module underlying $A$ is projective; it is therefore an
injective object of $\mathbf{A}$. Assertion~(b) follows from the fact
that $A$ is the ring of endomorphisms of this injective object.

Proposition~13 being proved, we turn our attention to the functor
$M\mapsto\operatorname{Hom}(M,I)$: the abelian group
$\operatorname{Hom}(M,I)$ is equipped with a left module structure over
the ring $A$ of endomorphisms of $I$. On the other hand, when $N$ runs
through the subobjects of finite length of $M$, the canonical
epimorphism from $M$ onto $M/N$ defines an injection of
$\operatorname{Hom}(M/N,I)$ into $\operatorname{Hom}(M,I)$. The image of
this map is a submodule which we shall again denote $l(N)$. In
general, the arguments we used when $M$ was equal to $I$ remain valid
here and they have the following consequence: the submodules $l(N)$ of
$\operatorname{Hom}(M,I)$ form a base of neighbourhoods of $0$ for a
topology which makes $\operatorname{Hom}(M,I)$ a pseudocompact
$A$-module. We shall always equip $\operatorname{Hom}(M,I)$ with this
topology, which we shall call \emph{natural}.

If $f\colon M\to M'$ is a morphism of $\mathbf{A}$,
$\operatorname{Hom}(f,I)$ is a continuous map from
$\operatorname{Hom}(M',I)$ into $\operatorname{Hom}(M,I)$: since
$\operatorname{Hom}(f,I)$ is a linear map, it suffices indeed to prove
that, for every subobject of finite length $N$ of $M$, there exists a
subobject of finite length $N'$ of $M'$ such that
$\operatorname{Hom}(f,I)$ maps $l(N')$ into $l(N)$. But it suffices to
choose $N'$ equal to $f(N)$. This shows that $\operatorname{Hom}(f,I)$
is a morphism of pseudocompact modules. We have therefore defined a
functor from $\mathbf{A}$ into the dual category of
$\operatorname{PC}(A)$.

\begin{theorem}
Let $I$ be an injective object of a locally finite category
$\mathbf{A}$, $A$ the ring of endomorphisms of $I$ and $F$ the functor
$M\mapsto\operatorname{Hom}(M,I)$. The functor $F$ defines by passing
to the quotient an equivalence between
$\mathbf{A}/\operatorname{Ker}F$ and the dual category of the category
$\operatorname{PC}(A)$ of pseudocompact modules over $A$.
\end{theorem}

The subcategory $\operatorname{Ker}F$ is formed of the $M$ such that
$\operatorname{Hom}(M,I)=0$ and it is localizing (corollary of
Proposition~8, Chap.~III). By Proposition~10 (Chap.~III),
$\operatorname{Ker}F$ is the smallest localizing subcategory containing
the simple objects $S$ such that $\operatorname{Hom}(S,I)$ is zero. One
may also say that the localizing subcategory $\operatorname{Ker}F$ is
generated by the simple objects which ``do not occur'' in the
decomposition of the socle of $I$ into isotypic components.

Let $T$ be the canonical functor from $\mathbf{A}$ into
$\mathbf{A}/\operatorname{Ker}F$. Then $TI$ is an injective object of
$\mathbf{A}/\operatorname{Ker}F$ and the socle of $TI$ contains simple
objects of $\mathbf{A}/\operatorname{Ker}F$ of every type. Moreover,
since $I$ is $\operatorname{Ker}F$-closed, the map $T(M,I)$ from
$\operatorname{Hom}(M,I)$ into $\operatorname{Hom}(TM,TI)$ is a
bijection. We may therefore replace $\mathbf{A}$ by the locally finite
category $\mathbf{A}/\operatorname{Ker}F$ and replace $I$ by $TI$. In
other words, we may reduce to the case where the category
$\operatorname{Ker}F$ is zero: this is what we shall suppose in the
sequel of this proof.

\begin{lemma}
For every object $M$ of $\mathbf{A}$, the map $F(M,I)$ from
$\operatorname{Hom}_{\mathbf{A}}(M,I)$ into
$\operatorname{Hom}_{\operatorname{PC}(A)}(FI,FM)$ is a bijection.
\end{lemma}

Let indeed $f$ be a morphism from $M$ into $I$. The image of $f$ under
the map $F(M,I)$ is none other than $\operatorname{Hom}_{\mathbf{A}}(f,I)$;
this latter map associates to every element $a$ of
$\operatorname{Hom}_{\mathbf{A}}(I,I)$ the element $a\circ f$ of
$\operatorname{Hom}_{\mathbf{A}}(M,I)$. But the morphisms from $A=FI$
into $FM$ are of the form $a\mapsto a\circ f$. This proves the lemma.

\begin{lemma}
Let $M$ be an object of $\mathbf{A}$, $(I_j)_{j\in E}$ a family of
objects isomorphic to $I$ and $J$ the direct sum of the family
$(I_j)_{j\in E}$. The map $F(M,J)$ from
$\operatorname{Hom}_{\mathbf{A}}(M,J)$ into
$\operatorname{Hom}_{\operatorname{PC}(A)}(FJ,FM)$ is bijective.
\end{lemma}

We first consider the case where the object $M$ is of finite length.
One knows that then the canonical map $u$ from the direct sum of the
groups $\operatorname{Hom}_{\mathbf{A}}(M,I_j)$ into
$\operatorname{Hom}_{\mathbf{A}}(M,J)$ is a bijection (Corollary~1 of
Theorem~1, Chap.~II); the same holds for the canonical map $v$ from the
direct sum of the $\operatorname{Hom}_{\operatorname{PC}(A)}(FI_j,FM)$
into $\operatorname{Hom}_{\operatorname{PC}(A)}(FJ,FM)$. Moreover,
$F(M,I_j)$ is a bijection for every $j$. The assertion therefore
follows from the equality
\[
v\circ\bigl(\sum F(M,I_j)\bigr)=F(M,J)\circ u.
\]

In the general case, $\operatorname{Hom}_{\mathbf{A}}(M,J)$ is the
projective limit of the groups $\operatorname{Hom}_{\mathbf{A}}(N,J)$
when $N$ runs through the subobjects of finite length of $M$.
Likewise, $\operatorname{Hom}_{\operatorname{PC}(A)}(FJ,FM)$ is the
projective limit of the groups
$\operatorname{Hom}_{\operatorname{PC}(A)}(FJ,FN)$. We have moreover the
following commutative diagram:
% Braced commas in the arrow labels, as required by tikz-cd.
\[
\begin{tikzcd}
\operatorname{Hom}_{\mathbf{A}}(M,J) \arrow[r, "F(M{,}J)"] \arrow[d, "\wr"'] & \operatorname{Hom}_{\operatorname{PC}(A)}(FJ,FM) \arrow[d, "\wr"'] \\
\varprojlim_N\operatorname{Hom}_{\mathbf{A}}(N,J) \arrow[r, "\varprojlim_N F(N{,}J)"] & \varprojlim_N\operatorname{Hom}_{\operatorname{PC}(A)}(FJ,FN)
\end{tikzcd}
\]

Since $F(N,J)$ is an isomorphism for every $N$, the same holds for
$F(M,J)$, which proves Lemma~2.

If the category $\operatorname{Ker}F$ is zero, let us denote by
$\mathbf{I}$ the full subcategory of $\mathbf{A}$ whose objects are
isomorphic to the direct sum of a family of objects equal to $I$.
Likewise, let $\mathbf{L}$ be the full subcategory of
$\operatorname{PC}(A)$ whose objects are isomorphic to the direct
product of a family of pseudocompact modules equal to $A$. Lemma~2
shows that the functor $F$ induces an equivalence between $\mathbf{I}$
and the category $\mathbf{L}^\circ$ dual to $\mathbf{L}$. Corollary~3
of Proposition~14 (Chap.~I) shows that $F$ defines an equivalence
between $\mathbf{A}$ and the category $\operatorname{PC}(A)^\circ$.

\begin{corollary}
With the hypotheses and notations of Theorem~4: if $I$ contains an
indecomposable injective of each type, the functor $F$ defines an
equivalence between $\mathbf{A}$ and the dual category of
$\operatorname{PC}(A)$.
\end{corollary}

We shall suppose from now on that $I$ contains an indecomposable
injective of each type. If $M$ is any subobject of $I$, we denote by
$l(M)$ the left ideal of $A$ formed of the endomorphisms of $I$ whose
kernel contains $M$.

\begin{corollary}
With the hypotheses and notations of Corollary~1: the map
$M\mapsto l(M)$ is a bijection from the set of subobjects of $I$ onto
the set of closed left ideals of $A$. The ideal $l(M)$ is two-sided if
and only if the subobject $M$ of $I$ is characteristic [i.e.\ $f(M)$
is contained in $M$ for every endomorphism $f$ of $I$].
\end{corollary}

The proof of Corollary~2 is left to the reader. We associate to every
closed subcategory $\mathbf{C}$ of $\mathbf{A}$ the largest subobject
$\mathbf{C}I$ of $I$ belonging to $\mathbf{C}$. If $f$ is an
endomorphism of $I$, $f(\mathbf{C}I)$ is contained in $\mathbf{C}I$. In
other words, $\mathbf{C}I$ is a characteristic subobject of $I$.

Conversely, let $J$ be a characteristic subobject of $I$. The objects
$M$ for which the canonical map from $\operatorname{Hom}(M,J)$ into
$\operatorname{Hom}(M,I)$ is bijective form a closed subcategory
$J\mathbf{A}$ of $\mathbf{A}$. Moreover, if $I$ contains an
indecomposable injective of each type, the maps
$\mathbf{C}\mapsto\mathbf{C}I$ and $J\mapsto J\mathbf{A}$ define a
bijective correspondence between the set of closed subcategories of
$\mathbf{A}$ and the set of characteristic subobjects of $I$. In other
words, the map $\mathbf{C}\mapsto l(\mathbf{C}I)$ is a bijection from
the set of closed subcategories of $\mathbf{A}$ onto the set of closed
two-sided ideals of $A$. In the corollary which follows, we write
$l(\mathbf{C})$ instead of $l(\mathbf{C}I)$.

\begin{corollary}
With the hypotheses and notations of Corollary~1: if $\mathbf{C}$ and
$\mathbf{D}$ are two closed subcategories of $\mathbf{A}$, the
two-sided ideal $l(\mathbf{C}\cdot\mathbf{D})$ is the closure of the
product $l(\mathbf{C})\cdot l(\mathbf{D})$.
\end{corollary}

The proof of Corollary~3 is left to the reader. It follows that with
the notations of Proposition~12, $\mathfrak{r}^{(n)}$ is the closure of
the product $\mathfrak{r}^n$.

If the injective $I$ contains one and only one indecomposable injective
of each type, the socle $I_0$ of $I$ contains one and only one simple
object of each type. One then says that the ring $A$ of endomorphisms
of $I$ is the \emph{pseudocompact ring associated to} $\mathbf{A}$.
This ring is determined up to isomorphism by the datum of the category
$\mathbf{A}$. Two equivalent locally finite categories have isomorphic
associated pseudocompact rings.

We say that a pseudocompact ring $A$ is \emph{sober} (\emph{sobre}) if
the quotient of $A$ by its Jacobson radical is a product of fields. We
say that two pseudocompact rings $A$ and $B$ are \emph{equivalent} if
the categories $\operatorname{PC}(A)$ and $\operatorname{PC}(B)$ are
equivalent. We then have the following results: the pseudocompact ring
associated to a locally finite category is sober. Every pseudocompact
ring is equivalent to a sober pseudocompact ring. Two sober
pseudocompact rings are equivalent if and only if they are isomorphic.

If $A$ is a pseudocompact ring, Theorem~4 applies in particular to the
category $\operatorname{dis}(A)$ of discrete topological $A$-modules.
If $A^*$ is the pseudocompact ring associated to
$\operatorname{dis}(A)$, the categories $\operatorname{dis}(A)$ and
$\operatorname{PC}(A^*)$ are dual [i.e.\ the categories
$\operatorname{dis}(A)$ and $\operatorname{PC}(A^*)^\circ$ are
equivalent]. It follows that the categories $\mathbf{T}(A)$ and
$\mathbf{T}(A^*)$, formed of the discrete modules of finite length over
$A$ and $A^*$, are dual. By Theorem~1 (Chap.~II) finally, every duality
between $\mathbf{T}(A)$ and $\mathbf{T}(A^*)$ extends to a duality
between $\operatorname{PC}(A)$ and $\operatorname{dis}(A^*)$. If $A$ is
sober, $A$ is therefore isomorphic to $(A^*)^*$ and one may say that
$A^*$ is the dual pseudocompact ring of $A$.

Corollary~1 has the following consequence:

\begin{corollary}
If $\mathbf{A}$ is a locally finite category and $A$ is the
pseudocompact ring associated to $\mathbf{A}$, then $\mathbf{A}$ is
equivalent to the category of discrete topological modules over the
ring $A^*$ dual to $A$.
\end{corollary}

\begin{corollary}
Let $I$ be an injective object of a locally finite category
$\mathbf{A}$ and $A$ the ring of endomorphisms of $I$. If $I$ contains
an indecomposable injective of each type, the centre $Z[\mathbf{A}]$ of
the category $\mathbf{A}$ (Chap.~III, \S~5) is isomorphic to the centre
of the ring $A$.
\end{corollary}

An element $z$ of $Z[\mathbf{A}]$ is indeed a morphism from the
identity functor $I_{\mathbf{A}}$ into itself. In particular, $z$
defines an endomorphism $z(I)\colon I\to I$. Since $z(I)$ commutes with
all the endomorphisms of $I$, $z(I)$ belongs to the centre of $A$. We
leave to the reader the task of verifying that the map $z\mapsto z(I)$
is a bijection from $Z[\mathbf{A}]$ onto the centre of $A$.

\begin{corollary}
Every commutative pseudocompact ring $A$ is isomorphic to the dual ring
$A^*$. In particular, the category $\mathbf{T}(A)$ of discrete
topological $A$-modules of finite length is equivalent to the dual
category $\mathbf{T}(A)^\circ$.
\end{corollary}

Since $A$ is a sober pseudocompact ring, it suffices indeed to show
that the pseudocompact rings $A$ and $A^*$ are equivalent, i.e.\ that
the categories $\mathbf{T}(A)$ and $\mathbf{T}(A)^\circ$ are equivalent.
We shall exhibit a contravariant functor from $\mathbf{T}(A)$ into
$\mathbf{T}(A)$ which defines an equivalence between $\mathbf{T}(A)$
and $\mathbf{T}(A)^\circ$.

When $\mathfrak{m}$ runs through the maximal open ideals of $A$, the
quotients $A/\mathfrak{m}$ run through simple objects of each type of
the category $\operatorname{dis}(A)$. We denote by $E$ an injective
envelope [in $\operatorname{dis}(A)$] of the direct sum of these
quotients.

For every $A$-module $M$, the abelian group $\operatorname{Hom}_A(M,E)$
is then equipped with an $A$-module structure: if $\varphi$ is an
element of $\operatorname{Hom}_A(M,E)$ and if $a$ is an element of $A$,
$a\cdot\varphi$ is the map $x\mapsto a\cdot\varphi(x)=\varphi(a\cdot x)$
from $M$ into $E$. In particular, the annihilator of
$\operatorname{Hom}_A(M,E)$ contains the annihilator of $M$. It follows
that $\operatorname{Hom}_A(M,E)$ is a discrete topological $A$-module
if $M$ is an object of $\mathbf{T}(A)$. Moreover,
$\operatorname{Hom}_A(M,E)$ is isomorphic to $A/\mathfrak{m}$ if $M$ is
isomorphic to $A/\mathfrak{m}$. The exactness of the functor
$M\mapsto\operatorname{Hom}_A(M,E)$ therefore implies that
$\operatorname{Hom}_A(M,E)$ is a discrete topological $A$-module of the
same length as $M$ when $M$ runs through the objects of
$\mathbf{T}(A)$. We denote by $D$ the contravariant functor from
$\mathbf{T}(A)$ into $\mathbf{T}(A)$ which we have just defined:
$DM=\operatorname{Hom}_A(M,E)$.

It suffices to show that the functor $D\circ D$ is isomorphic to the
identity functor of $\mathbf{T}(A)$ (Chap.~I, Proposition~12): if $M$
is a discrete topological $A$-module of finite length and $m$ is an
element of $M$, we denote by $m'$ the $A$-linear map from
$\operatorname{Hom}_A(M,E)$ into $E$ defined by the formula
$m'(f)=f(m)$. When $M$ varies, the maps $m\mapsto m'$ define a morphism
from the identity functor of $\mathbf{T}(A)$ into $D\circ D$. Moreover,
the map $m\mapsto m'$ is bijective when $M$ is of the form
$A/\mathfrak{m}$. The exactness of the functor $D\circ D$ implies that
it is still bijective when $M$ is of finite length. This completes the
proof.
